Este capítulo describe el proceso de construcción y puesta en marcha del \enquote{Asistente de Preparación \glslink{pmp}{PMP}}. Se detalla cómo el diseño arquitectónico propuesto en el capítulo anterior fue transformado en un sistema de software funcional y operativo. Se abordan las herramientas de desarrollo seleccionadas, la configuración del entorno de trabajo, la construcción de la interfaz de usuario, la integración de la \gls{ia} Generativa y la gestión segura de los datos. Asimismo, se discuten los desafíos encontrados durante la implementación y las soluciones adoptadas para garantizar un sistema robusto, seguro y fácil de usar.

\section{Entorno de desarrollo y estándares}

La fase de implementación comenzó con la definición de un entorno de trabajo estandarizado, diseñado para garantizar la calidad del código y facilitar el mantenimiento a largo plazo. Se priorizó el uso de herramientas líderes en la industria que permiten una gestión eficiente del ciclo de vida del software.

\subsection{Selección y justificación del stack tecnológico}
Para la construcción del sistema se seleccionó un conjunto de tecnologías modernas (\enquote{modern stack}) que favorecen el desarrollo rápido, la seguridad y el rendimiento. La elección de cada componente responde a necesidades específicas del proyecto, buscando siempre el equilibrio entre innovación y estabilidad.

\begin{table}[H]
\centering
\small
\setlength{\tabcolsep}{4pt}
\renewcommand{\arraystretch}{1.15}
\begin{tabularx}{\textwidth}{lllX}
\hline
\textbf{Componente} & \textbf{Tecnología} & \textbf{Rol en el Sistema} & \textbf{Beneficio Clave} \\
\hline
\textbf{Framework Web} & Next.js & Núcleo de la Aplicación & Permite crear aplicaciones web rápidas, seguras y optimizadas para motores de búsqueda. \\
\hline
\textbf{Interfaz Usuario} & React & Construcción Visual & Facilita la creación de interfaces interactivas y modulares que responden al instante. \\
\hline
\textbf{Lenguaje} & TypeScript & Lógica de Programación & Añade robustez al código, previniendo errores comunes antes de que ocurran. \\
\hline
\textbf{Estilo Visual} & Tailwind CSS & Diseño y Maquetación & Permite un diseño moderno y adaptable a cualquier dispositivo (móvil o escritorio). \\
\hline
\textbf{Inteligencia Artificial} & Google Gemini & Motor Cognitivo & Modelo avanzado capaz de razonar, explicar conceptos complejos y procesar gran cantidad de información. \\
\hline
\textbf{Orquestación IA} & LangChain & Conector de IA & Facilita la comunicación fluida y estructurada entre nuestra aplicación y el modelo de IA. \\
\hline
\textbf{Base de Datos} & PocketBase & Almacenamiento & Sistema ligero y seguro para guardar perfiles, progresos y conversaciones. \\
\hline
\end{tabularx}
\caption{Tecnologías clave del proyecto}
\label{tab:5-1}
\end{table}

\subsection{Configuración del entorno de trabajo}
Antes de iniciar la construcción, se establecieron las bases del entorno de desarrollo sobre Windows 11, asegurando que todas las herramientas estuvieran alineadas con los estándares profesionales:

\begin{enumerate}
\item \textbf{Plataforma de ejecución:} se configuró el entorno para soportar aplicaciones de alto rendimiento, garantizando compatibilidad con las últimas versiones de las librerías utilizadas.
\item \textbf{Gestión de código fuente:} se implementó un sistema de control de versiones utilizando \textbf{Git} y \textbf{GitHub}. Esto permite:
    \begin{itemize}
    \item Mantener un historial detallado de todos los cambios.
    \item Facilitar el trabajo colaborativo y la revisión de código.
    \item Asegurar que el código fuente esté respaldado en la nube.
    \end{itemize}
\item \textbf{Entorno de Desarrollo (IDE):} Se utilizó \textbf{Visual Studio Code}, configurado con extensiones de productividad que ayudan a mantener el estilo de código consistente y detectar problemas de calidad en tiempo real.
\end{enumerate}

\section{Construcción de la interfaz de usuario (frontend)}

La interfaz visual, \gls{ui}, es el punto de contacto entre el estudiante y el sistema. Su construcción se centró en reducir la carga cognitiva, permitiendo al usuario concentrarse únicamente en su estudio. Se implementó una arquitectura de \gls{spa}, lo que significa que la navegación es fluida y no requiere recargas constantes de la página.

\subsection{Diseño modular y experiencia de usuario}
El sistema se construyó utilizando componentes modulares, piezas de software independientes que se ensamblan para formar la aplicación completa, respondiendo a las especificaciones funcionales descritas en el Anexo \ref{anexo:especificacion}.

\subsubsection{A. Onboarding y personalización}
La primera impresión es crucial. Se desarrolló un módulo de bienvenida interactivo que guía al usuario en sus primeros pasos. Este componente:
\begin{enumerate}
\item Personaliza la experiencia saludando al usuario por su nombre.
\item Explica brevemente la metodología de estudio (Socrático, Simulaciones, etc.).
\item Configura el entorno inicial según las preferencias del estudiante.
\end{enumerate}

\subsubsection{B. Gamificación y motivación}
Para mantener el compromiso del estudiante, se implementaron elementos visuales de gamificación. Cuando un usuario completa un nivel o supera un examen, el sistema proporciona retroalimentación visual inmediata (como animaciones de celebración), reforzando positivamente el logro y fomentando la continuidad en el estudio.

\subsection{El chat inteligente: centro de interacción}
El componente de chat es el corazón del asistente. Su implementación técnica se enfocó en la velocidad y la claridad:
\begin{enumerate}
\item \textbf{Respuesta en tiempo real (streaming):} A diferencia de los sistemas tradicionales que hacen esperar al usuario, el chat muestra la respuesta palabra por palabra a medida que la IA la genera. Esto crea una sensación de conversación fluida y natural.
\item \textbf{Formato enriquecido:} El chat no solo muestra texto plano. Es capaz de renderizar tablas, listas y negritas, facilitando la lectura y comprensión de conceptos estructurados.
\end{enumerate}

\subsection{Simulador de examen}
El simulador recrea las condiciones del examen real \gls{pmp}. Técnicamente, funciona como una máquina de estados que gestiona el ciclo completo de la prueba:
\begin{itemize}
\item \textbf{Estado de carga:} prepara las preguntas.
\item \textbf{Estado de pregunta:} presenta el desafío y captura la respuesta.
\item \textbf{Estado de revisión:} permite marcar preguntas para revisar después.
\item \textbf{Estado de resultados:} calcula el puntaje final y muestra el análisis de rendimiento.
\end{itemize}
Las respuestas se guardan temporalmente en el dispositivo del usuario para evitar pérdidas de información si hay una interrupción breve de internet.

\section{Lógica del servidor e integración de \gls{ia} (backend)}

El \enquote{cerebro} del sistema reside en el servidor. Esta capa actúa como un intermediario seguro entre el usuario y los servicios de \gls{ia} de Google.

\subsection{Conexión con el modelo cognitivo (Gemini)}
Se integró el modelo \textbf{Gemini 3.0 Flash} debido a su capacidad superior para procesar grandes volúmenes de texto (como los estándares del \glslink{pmbok}{PMBOK}) y su velocidad de respuesta. Esta integración permite que el asistente no solo responda preguntas, sino que entienda el contexto completo de una situación de proyecto.

\subsection{Personalidades de la \gls{ia} (adaptabilidad del tutor)}
Una de las características más innovadoras implementadas es la capacidad del sistema para cambiar la \enquote{personalidad} de la \gls{ia} dinámicamente. Dependiendo del modo de estudio elegido por el usuario, el servidor instruye a la \gls{ia} para comportarse de manera diferente:

\begin{enumerate}
\item \textbf{Tutor Estándar:} ofrece explicaciones claras y directas.
\item \textbf{Simulación de Crisis:} adopta el rol de un interesado difícil o conflictivo.
\item \textbf{Tutor Socrático:} en lugar de dar respuestas, formula preguntas guía para que el estudiante deduzca la solución.
\item \textbf{Examinador Estricto:} realiza preguntas de prueba y evalúa las respuestas con rigor.
\item \textbf{Simplificador:} explica conceptos complejos usando analogías sencillas (explícamelo como a un niño).
\end{enumerate}

Esta lógica se ejecuta en el servidor, asegurando que las instrucciones pedagógicas (prompts) sean consistentes y no puedan ser alteradas por el usuario.

\section{Gestión de datos y seguridad}

La persistencia de la información es fundamental para el seguimiento del progreso del estudiante. Se implementó una base de datos robusta y ligera (\textbf{PocketBase}) que gestiona toda la información del sistema.

\subsection{Estructura de la Información}
Los datos se organizan en colecciones lógicas que permiten un acceso rápido y ordenado:
\begin{itemize}
\item \textbf{Usuarios:} almacena perfiles, preferencias y credenciales de acceso seguro.
\item \textbf{Progreso:} registra el avance en los niveles y las estadísticas de estudio.
\item \textbf{Historial de Conversaciones:} guarda todas las interacciones con el asistente para que el estudiante pueda repasarlas en el futuro.
\item \textbf{Resultados de Simulacros:} archiva los detalles de cada examen realizado para analizar la evolución del rendimiento.
\end{itemize}

\subsection{Seguridad y privacidad}
La seguridad se implementó siguiendo el principio de \enquote{mínimo privilegio}. Se configuraron reglas de acceso estrictas (\gls{rls}) que garantizan que cada estudiante solo pueda ver y modificar sus propios datos. Nadie puede acceder al historial o progreso de otro usuario, asegurando la total confidencialidad del proceso de aprendizaje.

Para facilitar la adopción del sistema por parte de los usuarios finales, se ha elaborado un manual detallado que describe el funcionamiento de cada una de las características implementadas (ver Anexo \ref{anexo:manual}).
